\chapter{Bases ortonormales y desarrollos de Fourier}\label{cap:nuevo15}\label{chap:31}
\index{base!ortonormal}\index{Fourier, coeficientes de}

Una base de álgebra lineal expresa cada vector mediante una combinación \emph{finita}. En análisis necesitamos permitir límites de combinaciones finitas. La ortogonalidad hace controlable ese límite: la norma de una cola se calcula sumando cuadrados de coeficientes. Las proyecciones del capítulo~\ref{cap:nuevo14} permiten interpretar cada suma finita como la mejor aproximación disponible en las direcciones elegidas. Desarrollaremos la lectura 31 de \cite{valencia}, incluyendo el argumento de existencia, la construcción explícita de Gram--Schmidt y el ejercicio de Müntz en la norma integral que realmente aparece en el enunciado.

\section{Bases algebraicas y bases ortonormales}
\label{lectura:31}
Una \emph{base de Hamel} es un conjunto linealmente independiente cuyo espacio generado algebraicamente es todo el espacio. Por ejemplo, los vectores coordenados $(e_n)$ no constituyen una base de Hamel de $\ell^{2}$: una combinación finita tiene soporte finito, mientras $(1/n)$ tiene infinitas coordenadas no nulas.

\begin{definicion}
Un conjunto es \emph{total} si el cierre de su espacio generado es todo el espacio. Una \emph{base ortonormal} de un espacio con producto interno es un conjunto ortonormal total. En el espacio cero, la base ortonormal es el conjunto vacío.
\end{definicion}
\index{base!de Hamel}\index{conjunto total}
La independencia lineal de un conjunto ortonormal se verifica tomando productos internos: si $\sum_{j=1}^{n}\alpha_je_j=0$, el producto con $e_k$ da $\alpha_k=0$. No conviene llamar simplemente «base» a cualquier conjunto independiente total: sin ortogonalidad, estas condiciones no garantizan un desarrollo en serie con coeficientes únicos. Una base de Schauder incorpora precisamente propiedades adicionales de convergencia y unicidad. Distinguiremos siempre las expresiones \emph{base de Hamel}, \emph{base de Schauder} y \emph{base ortonormal}.

\begin{teorema}\label{ch31:existencia}
Todo espacio de Hilbert tiene una base ortonormal.
\end{teorema}
\begin{proof}
Consideremos todos los subconjuntos ortonormales de $H$, incluido el vacío, ordenados por inclusión. Si $\mathcal C$ es una cadena, su unión es ortonormal: un vector de la unión pertenece a algún miembro y tiene norma uno; dos vectores de la unión pertenecen a dos miembros comparables de la cadena y, por tanto, ambos están en uno de ellos. Si son distintos, son ortogonales. La unión es una cota superior. El lema de Zorn proporciona un conjunto ortonormal maximal $E$.

Sea $M=\cl{\spanop E}$. Es un subespacio cerrado de Hilbert, por lo que tiene proyección ortogonal. Si $M\ne H$, elijamos $x\notin M$ y sea $r=x-P_Mx$. Entonces $r\ne0$ y $r\perp M$. El vector
\[
 e=\frac{r}{\norm r}
\]
tiene norma uno y es perpendicular a cada miembro de $E$. Por tanto $E\cup\{e\}$ es ortonormal, contradiciendo la maximalidad. La normalización es necesaria: el residuo $r$ no tiene, en general, norma uno. Así $M=H$ y $E$ es una base ortonormal. Para $H=\{0\}$, el mismo argumento da $E=\varnothing$.
\end{proof}

\section{Bases de Schauder: coordenadas mediante series}\label{ampl:schauder}
Una lista de vectores puede aproximar todo el espacio mediante combinaciones finitas sin proporcionar una receta coherente de coordenadas. En una expansión en serie no se permite escoger de nuevo todos los coeficientes cada vez que aumenta la precisión: las sumas parciales deben proceder de una misma lista de coeficientes. Esta es la exigencia adicional que formaliza Schauder.

La definición y los ejemplos de sucesiones de esta sección se apoyan en Maximenko \cite{max-schauder}. Desarrollamos las pruebas y añadimos ejemplos que separan independencia, densidad, convergencia y dependencia del orden. No se afirma que todo espacio separable tenga una base de Schauder; la separabilidad es aquí una condición necesaria, no un criterio suficiente.
\index{base!de Schauder}\index{sumas parciales!operador de}\index{base!condicional}

\begin{definicion}
Una sucesión $(b_n)_{n\ge1}$ en un espacio normado $V$ es una \emph{base de Schauder} si para cada $x\in V$ existe una única sucesión de escalares $(a_n)$ tal que
\[
 x=\lim_{N\to\infty}\sum_{n=1}^Na_nb_n
\]
en la norma de $V$. Se escriben $a_n=b_n^*(x)$ y
$S_Nx=\sum_{n=1}^Nb_n^*(x)b_n$. El orden de $(b_n)$ forma parte de la definición. En dimensión finita se utiliza la versión finita, que coincide con una base algebraica.
\end{definicion}

La convergencia significa $\norm{x-S_Nx}\to0$, no sólo convergencia coordenada a coordenada en algún modelo particular. Las aplicaciones $b_n^*$ y $S_N$ son lineales: las sumas de dos expansiones convergentes dan una expansión de la suma, y su unicidad obliga a sumar los coeficientes; el argumento para multiplicación por escalares es el mismo. No usaremos aquí un teorema general sobre continuidad de estos funcionales: en los ejemplos que la requieren obtendremos cotas explícitas.

\begin{proposicion}\label{ampl:schauder-independencia}
Una base de Schauder es linealmente independiente y total.
\end{proposicion}
\begin{proof}
Si $\sum_{n=1}^Na_nb_n=0$ con algún coeficiente no nulo, al añadir coeficientes cero después de $N$ obtenemos una expansión del vector cero distinta de la expansión con todos los coeficientes nulos. Esto contradice unicidad. Además, cada vector es límite de sus sumas parciales, que pertenecen al espacio generado algebraicamente por los $b_n$. Por ello el cierre de ese espacio generado es todo $V$.
\end{proof}

\begin{proposicion}[Bases coordenadas]\label{ampl:schauder-coordenadas}
Los vectores coordenados $e_n$ forman una base de Schauder de $\ell^p$ para $1\le p<\infty$ y de $c_0$ con su norma uniforme. Sus coeficientes son las coordenadas del vector. No forman una base de Schauder de $\ell^\infty$.
\end{proposicion}
\begin{proof}
Para $x\in\ell^p$, las sumas parciales $S_Nx=\sum_{n=1}^Nx_ne_n$ satisfacen
\[
 \norm{x-S_Nx}_p^p=\sum_{n>N}|x_n|^p\longrightarrow0.
\]
En $c_0$, $\norm{x-S_Nx}_\infty=\sup_{n>N}|x_n|\to0$ por la definición de convergencia a cero. Para probar unicidad, la evaluación en una coordenada $k$ es continua porque $|x_k-y_k|\le\norm{x-y}_p$ o $\norm{x-y}_\infty$. Una serie $\sum a_ne_n$ que converge a $x$ tiene, desde su suma parcial número $k$, coordenada $k$ igual a $a_k$, de modo que $a_k=x_k$.

En $\ell^\infty$, cualquier combinación finita de $e_n$ se anula a partir de un índice. Su distancia a la sucesión constante $\mathbf1$ es al menos $1$. Así la lista ni siquiera es total en $\ell^\infty$ y no puede ser base de Schauder.
\end{proof}

\begin{proposicion}[El caso ortonormal]\label{ampl:schauder-ortonormal}
Si una base ortonormal de un espacio con producto interno puede enumerarse como $(e_n)_{n\ge1}$, esa enumeración es una base de Schauder. Los coeficientes son $\ip{x}{e_n}$. Toda permutación de la enumeración tiene la misma propiedad.
\end{proposicion}
\begin{proof}
Sea $p_N=\sum_{n=1}^N\ip{x}{e_n}e_n$. El residuo $x-p_N$ es perpendicular a $e_1,\ldots,e_N$, de modo que $p_N$ es el mejor aproximante en su espacio generado. Dado $\eps>0$, la totalidad proporciona una combinación finita $v$ de los $e_n$ con $\norm{x-v}<\eps$. Para $N$ suficientemente grande, $v$ pertenece al espacio generado por $e_1,\ldots,e_N$; entonces $\norm{x-p_N}\le\norm{x-v}<\eps$. Esto prueba convergencia sin suponer completitud adicional del ambiente.

Si $x=\lim_N\sum_{n=1}^Na_ne_n$, tomar producto interno con $e_k$ y usar su continuidad da $\ip{x}{e_k}=a_k$. La expansión es única. Después de una permutación, toda combinación finita sigue estando contenida en alguna suma inicial suficientemente larga, y el mismo argumento funciona. Para familias ortonormales no numerables se usan sumas sobre conjuntos finitos; la definición secuencial de Schauder no debe aplicarse sin esta distinción.
\end{proof}

\section{Ejercicios sobre coordenadas, densidad y orden}

\begin{ejercicio}[Ampliación: una base para las sucesiones convergentes]\label{ej:ampl-schauder-1}
En el espacio $c$ con la norma uniforme pruebe que $(\mathbf1,e_1,e_2,\ldots)$ es una base de Schauder. Determine los coeficientes, el error de las sumas parciales y cotas para los funcionales coordenados. Este ejercicio adapta la propuesta de encontrar una base de $c$ en \cite{max-schauder}; la construcción y la solución se desarrollan aquí.
\end{ejercicio}


\begin{ejercicio}[Ampliación: independencia y totalidad no bastan]\label{ej:ampl-schauder-2}
En $\ell^2$ defina $b_1=e_1$ y $b_n=e_1+n^{-1}e_n$ para $n\ge2$. Pruebe que la lista es linealmente independiente y total, pero no es una base de Schauder. Examine el vector $x=(0,1/2,1/3,\ldots)$. Ejercicio propio de esta edición.
\end{ejercicio}


\begin{ejercicio}[Ampliación: separabilidad y una obstrucción]\label{ej:ampl-schauder-3}
Pruebe que todo espacio con base de Schauder es separable. Demuestre después directamente que $\ell^\infty$ no es separable y concluya que no admite ninguna base de Schauder numerable, no sólo que su lista coordenada falla. La primera implicación y esta aplicación desarrollan la discusión de \cite{max-schauder}.
\end{ejercicio}


\begin{ejercicio}[Ampliación: transportar una base a otra]\label{ej:ampl-schauder-4}
Sea $T:V\to W$ una biyección lineal con $\norm{Tx}\le C\norm x$ y $\norm{T^{-1}y}\le D\norm y$. Demuestre que una base de Schauder $(b_n)$ de $V$ se transforma en la base $(Tb_n)$ de $W$. Aplique el resultado en $\ell^2$ a $T=I+\lambda R$, donde $R(x_1,x_2,\ldots)=(0,x_1,x_2,\ldots)$ y $|\lambda|<1$. Construya la inversa y calcule los coeficientes de la base $e_n+\lambda e_{n+1}$. Ejercicio propio de esta edición.
\end{ejercicio}


\begin{ejercicio}[Ampliación: una base cuyo orden importa]\label{ej:ampl-schauder-5}
En $c_0$ con norma uniforme, sea $s_n=e_1+\cdots+e_n$. Pruebe que $(s_n)$ es una base de Schauder, calcule sus coeficientes y la norma del operador de suma parcial. Usando $x_n=(-1)^{n-1}/n$, construya una permutación de esta base que deje de ser base de Schauder. Ejercicio propio de esta edición.
\end{ejercicio}


Hamel describe representaciones finitas; Schauder añade límites ordenados de una expansión única; una base ortonormal permite además medir el error con cuadrados de coeficientes. Desarrollaremos ahora esta última ventaja, que conduce a Parseval y a Riesz--Fischer.

\section{Coordenadas y teorema de Riesz--Fischer}
\begin{teorema}[Parseval y expansión]\label{ch31:parseval}
\index{Parseval, identidad de}
Si $E=(e_s)_{s\in S}$ es una base ortonormal de $H$, entonces
\[
 \norm{x}^{2}=\sum_{s\in S}|\ip{x}{e_s}|^{2},\qquad
 x=\sum_{s\in S}\ip{x}{e_s}e_s.
\]
La segunda suma converge en norma: para todo $\eps>0$ existe $F_0\subset S$ finito tal que cualquier suma sobre un conjunto finito $F\supset F_0$ dista menos de $\eps$ de $x$.
\end{teorema}
\begin{proof}
La identidad de Parseval del ejercicio~\ref{ej:lectura-30-2}, aplicada a $M=H$, da la primera fórmula. Para $p_F=\sum_{s\in F}\ip{x}{e_s}e_s$, la ortogonalidad da
\[
 \norm{x-p_F}^{2}=\norm{x}^{2}-\sum_{s\in F}|\ip{x}{e_s}|^{2}.
\]
Como las sumas finitas tienen supremo $\norm{x}^{2}$, existe $F_0$ para el cual la diferencia es menor que $\eps^{2}$. Al ampliar $F_0$ la suma solo aumenta, de modo que la misma estimación vale para todo $F\supset F_0$. Además el soporte de los coeficientes es numerable por Bessel; si se enumera, se obtiene una serie ordinaria con el mismo límite, independiente del orden.
\end{proof}

\begin{teorema}[Riesz--Fischer]\label{ch31:rieszfischer}
\index{Riesz--Fischer, teorema de}
La aplicación
\[
 U:H\longrightarrow\ell^{2}(S),\qquad Ux=(\ip{x}{e_s})_{s\in S},
\]
es una biyección lineal que preserva normas y productos internos.
\end{teorema}
\begin{proof}
Bessel garantiza que $Ux\in\ell^{2}(S)$. La linealidad en la primera entrada prueba que $U$ es lineal, y Parseval prueba $\norm{Ux}_2=\norm{x}$. En particular, es inyectiva.

Para la suprayectividad, sea $a\in\ell^{2}(S)$. Si su soporte es finito, $x=\sum_s a_se_s$ tiene las coordenadas requeridas. Si es infinito, enumerémoslo como $s_1,s_2,\dots$ y definamos $x_N=\sum_{j=1}^{N}a_{s_j}e_{s_j}$. Para $M>N$, Pitágoras da
\[
 \norm{x_M-x_N}^{2}=\sum_{j=N+1}^{M}|a_{s_j}|^{2}\longrightarrow0
\]
cuando $N\to\infty$, uniformemente para $M>N$. Por completitud existe $x=\lim_Nx_N\in H$. Para cada $s=s_j$, $\ip{x_N}{e_s}=a_s$ una vez que $N\geq j$; la continuidad del producto interno da $\ip{x}{e_s}=a_s$. Si $s$ está fuera del soporte, todos esos productos son cero y también lo es su límite. Luego $Ux=a$.

Finalmente, si $x_F,y_F$ son truncamientos sobre un conjunto finito común, $\ip{x_F}{y_F}=\sum_{s\in F}\ip{x}{e_s}\overline{\ip{y}{e_s}}$. Las aproximaciones en norma y la continuidad del producto interno permiten pasar al límite. La serie del lado derecho converge absolutamente por Cauchy--Schwarz en $\ell^{2}(S)$. Por ello $\ip{Ux}{Uy}_{\ell^{2}(S)}=\ip{x}{y}$.
\end{proof}
Los escalares $\ip{x}{e_s}$ se llaman \emph{coeficientes de Fourier} respecto de $E$. Esta denominación no exige funciones trigonométricas: describe cualquier sistema de coordenadas ortonormales.

\section{Ejercicios de ortonormalización, aproximación y Fourier}

\begin{ejercicio}[Lectura 31, ejercicio 1: Gram--Schmidt]\label{ej:lectura-31-1}
\index{Gram--Schmidt, procedimiento de}
Para una familia finita o numerable linealmente independiente $v_1,v_2,\dots$ en un espacio con producto interno, definir
\[
 w_1=v_1,\qquad
 w_{n+1}=v_{n+1}-\sum_{k=1}^{n}\frac{\ip{v_{n+1}}{w_k}}{\norm{w_k}^{2}}w_k.
\]
Probar la ortogonalidad, la no nulidad, la igualdad de los espacios generados y la existencia de una base ortonormal en todo espacio prehilbertiano separable.
\end{ejercicio}


\begin{ejercicio}[Lectura 31, ejercicio 2: primeros polinomios de Legendre]\label{ej:lectura-31-2}
\index{Legendre, polinomios de}
Aplicar Gram--Schmidt a $1,x,x^{2},x^{3},x^{4}$ con $\ip{f}{g}=\int_{-1}^{1}f(x)\overline{g(x)}\,dx$.
\end{ejercicio}


\begin{ejercicio}[Lectura 31, ejercicio 3]\label{ej:lectura-31-3}
Probar que $(e^{inx})_{n\in\Z}$ es ortonormal en $C([-\pi,\pi])$ con el producto interno normalizado.
\end{ejercicio}


\begin{ejercicio}[Lectura 31, ejercicio 4: Müntz en norma integral]\label{ej:lectura-31-4}
\index{Muntz@Müntz, teorema de}\index{Gram, determinante de}
Sean $1\leq n_1<n_2<\cdots$ enteros y $M_k=\spanop\{x^{n_1},\dots,x^{n_k}\}\subset C([0,1])$, con $\norm{f}_2^{2}=\int_0^1|f|^{2}$. Probar que $\bigcup_kM_k$ es denso para esta norma si y solo si $\sum_k1/n_k=\infty$. Desarrollar (i) el cociente de determinantes de Gram y la fórmula de distancia; (ii) el criterio de las series logarítmicas con sus hipótesis; (iii) el análisis del producto resultante.
\end{ejercicio}


\begin{ejercicio}[Lectura 31, ejercicio 5]\label{ej:lectura-31-5}
Probar que un espacio de Hilbert es separable si y solo si tiene una base ortonormal a lo sumo numerable.
\end{ejercicio}


\begin{ejercicio}[Lectura 31, ejercicio 6]\label{ej:lectura-31-6}
Calcular los cinco primeros coeficientes respecto de $e_0,\dots,e_4$ del ejercicio~\ref{ej:lectura-31-2} para (i) $x^2+x+1$, (ii) $\sin(\pi x)$, (iii) $\cos(\pi x)$ y (iv) $|x|$.
\end{ejercicio}


\begin{ejercicio}[Lectura 31, ejercicio 7]\label{ej:lectura-31-7}
Si $E$ es ortonormal y $x\in H$, estimar el número $n_k$ de vectores $e\in E$ con $|\ip{x}{e}|>1/k$.
\end{ejercicio}


\section{Una base como sistema de medición}
Los coeficientes de Fourier registran las componentes detectadas por cada dirección ortonormal. Bessel limita la energía registrada; Parseval afirma que una base total no pierde ninguna componente; Riesz--Fischer garantiza que toda lista de coeficientes de energía finita procede de un vector. La completitud permite pasar de las sumas finitas, puramente algebraicas, al vector que esas sumas aproximan. En el capítulo~\ref{cap:nuevo16} pasaremos de las coordenadas de un vector a la representación de los funcionales continuos y a la relación de adjunción entre operadores.
